%! Author = sriadityavedantam
%! Date = 3/30/26

\section{Group Actions}

Let $G$ be a finite group and let $A$ be a set.
We say $G$ acts on $A$ if there is a homomorphism
\[
    G\to \sym(A).
\]
\begin{definition}[Left Translation Group Action]
    A left group action of $G$ on $A$ is a rule assigning to each $g\in G$ and $a\in A$ an element
    \[
        g\cdot a\in A
    \]
    such that
    \[
        1\cdot a=a
    \]
    and
    \[
        (gh)\cdot a=g\cdot (h\cdot a).
    \]
    The most general example is the action of $G$ on itself by left translation, also called the Cayley action, where
    \[
        g\mapsto \phi_g,
        \qquad
        \phi_g(x)=gx.
    \]
\end{definition}

\begin{definition}[Right Conjugation Action]
    A right action of $G$ on $A$ is a rule assigning to each $g\in G$ and $a\in A$ an element
    \[
        a^g\in A
    \]
    such that
    \[
        a^1=a
    \]
    and
    \[
        a^{gh}=(a^g)^h.
    \]
\end{definition}

\begin{definition}[Left Conjugation Action]
    A left action of $G$ on $A$ may also be written as
    \[
            {}^{g}a
    \]
    for $g\in G$ and $a\in A$, such that
    \[
            {}^{1}a=a
    \]
    and
    \[
            {}^{gh}a = {}^g({}^{h}a).
    \]
\end{definition}

\begin{definition}[Orbit]
    Suppose $G$ acts on $A$ by left translation.
    Let $a\in A$.
    Then the orbit of $a$ under the action is
    \[
        \{g\cdot a:g\in G\}\eqqcolon \mathcal{O}(a).
    \]
\end{definition}

\begin{lemma}{
    For the Cayley action of $G$ on itself, the orbit of any element is all of $G$.
}
    Let $a\in G$.
    Then
    \[
        \mathcal{O}(a)=\{ga:g\in G\}.
    \]
    Given any $x\in G$, choose
    \[
        g=xa^{-1}.
    \]
    Then
    \[
        g\cdot a=(xa^{-1})a=x.
    \]
    So every element of $G$ lies in the orbit of $a$.
    Hence
    \[
        \mathcal{O}(a)=G.
    \]
\end{lemma}

If the action is written on the right, then
\[
    \mathcal{O}(a)=\{a^g:g\in G\}.
\]
If $G$ acts on itself by conjugation, then
\[
    \mathcal{O}(x)=\{gxg^{-1}:g\in G\},
\]
which is the conjugacy class of $x$.
We denote this by $\mathcal{C}_x$.
If $G$ is abelian, then
\[
    \mathcal{C}_x=\{x\}
\]
for every $x\in G$.

\begin{example}
    If
    \[
        G=S_3,
    \]
    then
    \[
        \mathcal{C}_{(12)}=\{(12),(13),(23)\}.
    \]
    Also,
    \[
        G=\mathcal{C}_1\sqcup \mathcal{C}_{(12)}\sqcup \mathcal{C}_{(123)}.
    \]
\end{example}

If $G$ acts on $A$, then we can define a relation on $A$ by declaring
\[
    a\sim b
\]
if $a$ and $b$ are in the same orbit.
This is an equivalence relation.

\begin{definition}[Stabilizer]
    Suppose $G$ acts on $A$.
    Let $a\in A$.
    Then the stabilizer of $a$ under the action is
    \[
        \stab_G(a)\coloneqq \{g\in G:g\cdot a=a\}
    \]
    for a left action, and
    \[
        \stab_G(a)\coloneqq \{g\in G:a^g=a\}
    \]
    for a right action.
\end{definition}

If $G$ acts on itself by conjugation, then the stabilizer of $x$ is exactly the centralizer
\[
    C_G(x).
\]

\begin{theorem}[Orbit-Stabilizer Theorem]{
    Suppose $G$ acts on a set $A$ and let $a\in A$.
    Then:
    \begin{enumerate}
        \item $\stab_G(a)\leq G$;
        \item
        \[
            [G:\stab_G(a)] = \ord(\mathcal{O}(a)).
        \]
    \end{enumerate}
}
    \textbf{(1)}
    Let $H=\stab_G(a)$.
    Since
    \[
        1\cdot a=a,
    \]
    we have
    \[
        1\in H.
    \]
    If $g,h\in H$, then
    \[
        h\cdot a=a.
    \]
    So
    \[
        (gh)\cdot a=g\cdot(h\cdot a)=g\cdot a=a.
    \]
    Hence
    \[
        gh\in H.
    \]
    Now if $g\in H$, then
    \[
        g\cdot a=a.
    \]
    Apply $g^{-1}$ to both sides:
    \[
        g^{-1}\cdot(g\cdot a)=g^{-1}\cdot a.
    \]
    Thus
    \[
        a=g^{-1}\cdot a,
    \]
    so
    \[
        g^{-1}\in H.
    \]
    Therefore
    \[
        H\leq G.
    \]

    \textbf{(2)}
    We find a bijection between the set of left cosets of $H$ in $G$ and the orbit of $a$.
    Define
    \[
        gH \longmapsto g\cdot a.
    \]
    We need to check that this is well-defined.
    Now
    \[
        g_{1}H=g_{2}H
        \iff g_2^{-1}g_1\in H
        \iff (g_2^{-1}g_1)\cdot a=a.
    \]
    This is equivalent to
    \[
        g_1\cdot a=g_2\cdot a.
    \]
    So the map is well-defined and injective.
    It is surjective by definition of the orbit.
    Hence
    \[
        [G:H]=\ord(\mathcal{O}(a)).
    \]
    That is,
    \[
        [G:\stab_G(a)] = \ord(\mathcal{O}(a)).
    \]
\end{theorem}

\begin{definition}[Class Equation]
    Let $G$ act on itself by conjugation.
    Then $G$ is the union of its disjoint conjugacy classes.
    Suppose these have representatives
    \[
        z_1=1,z_2,\ldots,z_k,x_1,x_2,\ldots,x_r,
    \]
    where the $z_i$ are the elements of the center $Z(G)$, so their conjugacy classes have size $1$.
    Then the class equation for $G$ is
    \[
        \ord(G)=\ord(Z(G))+\ord(\mathcal{C}_{x_1})+\cdots+\ord(\mathcal{C}_{x_r}).
    \]
    Using orbit-stabilizer, this may also be written as
    \[
        \ord(G)=\ord(Z(G))+\sum_{i=1}^r [G:C_G(x_i)].
    \]
\end{definition}

\begin{example}
    Suppose $P$ is a group with
    \[
        \ord(P)=p^n,
    \]
    where $p$ is prime.
    Then
    \[
        Z(P)\neq \{1\}.
    \]
\end{example}

\begin{proof}
    From the class equation,
    \[
        \ord(P)=\ord(Z(P))+\sum_{i=1}^r [P:C_P(x_i)].
    \]
    Now for each $x_i\notin Z(P)$,
    \[
        [P:C_P(x_i)] = \frac{p^n}{\ord(C_P(x_i))}.
    \]
    Since $C_P(x_i)$ is a proper subgroup of $P$, its order is $p^m$ for some $m<n$.
    Thus
    \[
        [P:C_P(x_i)]=p^{n-m},
    \]
    which is divisible by $p$.
    Therefore every term in the sum
    \[
        \sum_{i=1}^r [P:C_P(x_i)]
    \]
    is divisible by $p$.
    Since $\ord(P)=p^n$ is also divisible by $p$, it follows that
    \[
        \ord(Z(P))
    \]
    is divisible by $p$.
    Hence
    \[
        \ord(Z(P))\neq 1.
    \]
    So
    \[
        Z(P)\neq \{1\}.
    \]
\end{proof}